% ===================================================================
%  TÁBLA Uimh. 11 — An Chathair agus a Cuid Foirgneamh. — An Dóiteán.
%  Traduction 2026 d'apres texto/io, controlee sur texto/fr, texto/en
%  et texto/es. Consigne : voir texto/ga/10-tabla-01.tex.
%
%  OUVERTURE DE LA TROISIEME SERIE : « AN TRÍÚ SRAITH ». Comme apres
%  « an dara », « sraith » ne mute pas apres « an tríú » — voir la
%  note du tableau 7 sur « sh?raith ».
%
%  LE BLOC c1-03-1 A COUTE UNE INVERSION AU FINNOIS ET UNE AU BASQUE :
%  « la pordo (8) dil olda fortifikita kastelo (7) ». L'irlandais
%  postpose son genitif comme l'ido — « geata (8) an tseanchaisleáin
%  dhaingnithe (7) » — et ne paie rien. C'est la onzieme fois de cette
%  colonne que la source ordinaire des inversions ne joue pas.
%
%  MAIS LE BLOC c2-02-1 EN A COUTE UNE, et elle vient encore du VERBE.
%  L'ido ecrit « plu bone kam la soldati (57) o la jendarmi (63),
%  retroirigas la kuriozi (56) » : le terme de comparaison AVANT
%  l'objet. L'irlandais range verbe, sujet, objet, et rejette la
%  comparaison a la fin — l'objet (56) passait donc devant (57) et
%  (63). Le remede n'est pas de deplacer un renvoi mais de retourner
%  la phrase : on la met au negatif et l'on prend les soldats pour
%  sujet — « Ní chuireann na saighdiúirí (57) ná na gardaí (63) na
%  fiosraitheoirí (56) siar chomh maith leis na marcaigh ». L'ordre de
%  l'ido revient, et la phrase dit la meme chose.
% ===================================================================

%%K t11-apar-1 apar
\VUsaut{2.0mm}
\VUcentre{13.2pt}{90}{AN TRÍÚ SRAITH}
\VUsaut{3.0mm}
\VUfilet{16mm}
\VUsaut{5.6mm}
\VUtitre{15.0pt}{60}{TÁBLA Uimh. 11}
\VUsaut{1.4mm}
\VUpk{11.4pt}{40}{An Chathair agus a Cuid Foirgneamh. --- An Dóiteán.}
\VUsaut{2.2mm}

%%K t11-tit-1 sub
\VUpk{10.2pt}{40}{An Bruachbhaile.}
\VUsaut{3.0mm}

%%K t11-c1-01-1 p
1. --- Tá cónaí orm i gcathair mhór atá céad ciliméadar ón Aigéan, agus is
\VUgras{calafort} den chéad ghrád í ina dhiaidh sin \textsuperscript{(1)}. Tá an
abhainn a shníonn lena taobh ceithre chéad méadar ar leithead, agus déanann sí ród
an-bhreá.

%%K t11-c1-02-1 p
2. --- An chuid sin uile den chathair atá ar na \VUgras{céanna}
\textsuperscript{(2)} tá cuma mhara agus tionsclaíoch ar fad uirthi, agus déanann sí
\VUgras{bruachbhaile} \textsuperscript{(3)} nach gcónaíonn ann ach mairnéalaigh agus
lucht oibre. Oibríonn na daoine sin sna \VUgras{monarchana} \textsuperscript{(4)} agus
sna \VUgras{gasoibreacha} \textsuperscript{(5)} a bhfeictear a \VUgras{simléar} ard
\textsuperscript{(6)} agus a \VUgras{ngasaiméadar} i bhfad uathu.

%%K t11-c1-03-1 p
3. --- Sa mheánaois bhí an chathair daingnithe, agus tá rian dá ballaí le fáil fós,
go háirithe \VUgras{geata} \textsuperscript{(8)} an \VUgras{tseanchaisleáin
dhaingnithe} \textsuperscript{(7)} agus a dhá \VUgras{thúr} \textsuperscript{(9)} ar
gach taobh de, túir a úsáidtear inniu mar \VUgras{phríosún}.

%%K t11-c1-04-1 p
4. --- Sa \VUgras{cheantar} sin ar fad, a bhfuil cuma an-ársa air, níl ach aon
fhoirgneamh amháin sách nua-aimseartha: sin \VUgras{teach an chustaim}
\textsuperscript{(10)}. Tá sé idir an ché agus an \VUgras{sráidín}
\textsuperscript{(11)} a théann chuig sean-\VUgras{halla na cathrach}
\textsuperscript{(12)}. Aithnítear an foirgneamh deireanach sin gan dua ar an
\VUgras{gclogás} ollmhór \textsuperscript{(13)} atá os cionn na seanchathrach.

%%K t11-c1-05-1 p
5. --- Ar an taobh eile den tsráid \VUgras{seasann} foirgneamh a tógadh ar an stíl
chéanna le teach an chustaim: sin an \VUgras{Mhalartlann} \textsuperscript{(14)}.
Féachann ceann dá aghaidheanna ar chearnóg bheag ina bhfuil \VUgras{teach an
phréifeachta} \textsuperscript{(15)}. Go deimhin, gan a bheith ina príomhchathair, is
príomhbhaile tábhachtach rannáin í ár gcathair.

%%K t11-tit-2 sub
\VUsaut{5.0mm}
\VUpk{10.2pt}{40}{An Chuid is Airde den Chathair.}
\VUsaut{3.4mm}

%%K t11-c1-06-1 p
6. --- Tá an garastún, atá sách líonmhar, ar lóistín i \VUgras{mbeairicí}
\textsuperscript{(16)} fairsinge an-fholláin; sínid chomh fada le ráillí \VUgras{pháirc
na cathrach} \textsuperscript{(17)}. Gabhann an \VUgras{bulbhard}
\textsuperscript{(19)} a théann chuig an bpáirc sin taobh thiar den \VUgras{bhanc
taisce} \textsuperscript{(18)}, agus baineann sé amach cearnóg na \VUgras{hardeaglaise}
\textsuperscript{(20)}.

%%K t11-c1-07-1 p
7. --- Tá na foirgnimh sin uile leata mar a bheadh amfaitéatar ar chnocán ina bhfuil
ár \VUgras{n-ardscoil} fhairsing \textsuperscript{(21)} chomh maith.

%%K t11-c1-07-2 p
Má théitear síos bóthar atá beagán ar fhána, bóthar a cheanglaíonn leis an tSráid
Mhór, tagtar chuig \VUgras{teach na cúirte} \textsuperscript{(22)}, a bhfuil
\VUgras{gairdín poiblí} taitneamhach \textsuperscript{(26)} sínte os a chomhair. Níl an
\VUgras{chearnóg} sin an-mhór, ach cuireann sí ósais glaise agus fionnuaire i lár na
cathrach. Dá bhrí sin is minic muintir na cathrach ann, iad tugtha do shuí faoi
phuball an \VUgras{chaife} \textsuperscript{(25)} chun éisteacht le ceoltóirí an airm a
sheinneann Déardaoin agus Domhnach sa \VUgras{bhothán ceoil} \textsuperscript{(27)}
atá curtha idir na faichí féir agus na ceapacha bláth.

%%K t11-c1-08-1 p
8. --- Ar cheann de na faichí sin seasann \VUgras{dealbh} chré-umha
\textsuperscript{(28)}, leacht a tógadh i gcuimhne ar dhuine dár gcomhchathróirí a
rinne seirbhísí móra dá chathair dhúchais.

%%K t11-tit-3 sub
\VUsaut{4.0mm}
\VUpk{10.2pt}{40}{Na Modhanna Iompair.}
\VUsaut{3.4mm}

%%K t11-c1-09-1 p
9. --- Go dtí seo níl ár gcathair soilsithe leis an solas leictreach, níl aici ach
\VUgras{lampaí gáis} \textsuperscript{(29)}. Ach tá \VUgras{na nuachtáin áitiúla} ag
feachtasaíocht chun go bhfeabhsófaí an córas soilsithe sin, \VUgras{amhail} mar a
feabhsaíodh cheana \VUgras{córas tarraingthe na dtramanna}. Cuireadh
\VUgras{tramanna leictreacha} \textsuperscript{(31)} beagnach in ionad na
sean-\VUgras{charranna} a bhíodh á \VUgras{dtarraingt} ag \VUgras{capaill}, agus
iompraíonn siad na taistealaithe go mear ó \VUgras{cheann na cathrach go dtí an
ceann eile} ar phraghas an-\VUgras{saor}.

%%K t11-c1-10-1 p
10. --- In aice le teach na cúirte tá an \VUgras{Banc} \textsuperscript{(32)} ina
lascainítear na sócmhainní tráchtála. Téann \VUgras{teachtairí an bhainc}
\textsuperscript{(33)} chun na dtithe chun na fiacha a bhailiú.

%%K t11-tit-4 sub
\VUsaut{4.0mm}
\VUpk{10.2pt}{40}{Ealaín agus Oideachas.}
\VUsaut{3.4mm}

%%K t11-c1-11-1 p
11. --- Is é an foirgneamh breá atá in aice leis an \VUgras{músaem}. Tá ann
\VUgras{leabharlann} \textsuperscript{(34)} na cathrach, \VUgras{bailiúcháin staire
nádúrtha} \textsuperscript{(35)} áille (an Músaem Nádúir) agus \VUgras{dánlann}
mhaisiúil \textsuperscript{(36)} atá ina \VUgras{taispeántas} buan den scoth.

%%K t11-c1-11-2 p
Osclaítear don phobal é dhá uair sa tseachtain, ach ceadaítear do choimhthígh cuairt a
thabhairt air lá ar bith ar shíneadh láimhe a thabhairt don \VUgras{choimeádaí}
\textsuperscript{(37)}. Tagann na \VUgras{péintéirí} \textsuperscript{(38)} ann chun
oibre. Feictear iad \VUgras{os comhair} a gcreata, agus \VUgras{pailéad} ina lámh, ag
cóipeáil na bpictiúr cáiliúil.

%%K t11-c1-12-1 p
12. --- Ar an ard, taobh thiar den mhúsaem, feictear \VUgras{cruinneachán}
\textsuperscript{(40)} an \VUgras{ospidéil} \textsuperscript{(39)} agus foirgnimh an
\VUgras{choláiste oiliúna} \textsuperscript{(41)} inar oileadh múinteoirí ár
scoileanna pobail. Ar chearnóg na hAmharclainne tá \VUgras{oifig phoist}
\textsuperscript{(43)} curtha i \VUgras{dteach príobháideach} \textsuperscript{(42)}.

%%K t11-tit-5 sub
\VUsaut{5.0mm}
\VUpk{11.4pt}{40}{An Dóiteán.}
\VUsaut{3.0mm}

%%K t11-tit-6 sub
\VUpk{10.2pt}{40}{Gairm Dhóiteáin.}
\VUsaut{3.4mm}

%%K t11-c2-01-1 p
1. --- Scaipeann ráfla scanrúil ar fud na cathrach: tá dóiteán tar éis briseadh amach
san amharclann! An nóiméad a shroichim an \VUgras{chearnóg} \textsuperscript{(96)}, tá
an slua cruinnithe cheana ar an \VUgras{gcosán} \textsuperscript{(46)} agus ar an
\VUgras{mbóthar} \textsuperscript{(47)}. Tagann \VUgras{oibrithe an phoist}
\textsuperscript{(44)} agus \VUgras{fir an phoist} \textsuperscript{(45)} amach as an
oifig. B'éigean do \VUgras{thiománaí tram} \textsuperscript{(48)} a charr a stopadh
chun \VUgras{otharcharr} na cathrach \textsuperscript{(50)} a ligean thart, carr a
thagann le \VUgras{máinlia} \textsuperscript{(52)} agus le \VUgras{haltra}
\textsuperscript{(51)} chun aire a thabhairt do \VUgras{dhuine gortaithe}
\textsuperscript{(53)} a tugadh ar \VUgras{shínteán} \textsuperscript{(54)} in aice
leis an \VUgras{mboth nuachtán} \textsuperscript{(55)}.

%%K t11-c2-02-1 p
2. --- Stad complacht coisithe, a bhí ag filleadh ón gcleachtadh, chun an t-ord a
choimeád i dteannta na \VUgras{ngardaí cathrach ar muin capaill}
\textsuperscript{(61)}. Ní chuireann na \VUgras{saighdiúirí} \textsuperscript{(57)} ná
na \VUgras{gardaí} \textsuperscript{(63)} na \VUgras{fiosraitheoirí}
\textsuperscript{(56)} siar chomh maith leis na \VUgras{marcaigh}, a bhfuil a
gcapaill ag éirí in airde ar a gcosa deiridh. Tá na gardaí an-ghnóthach mar sin féin.
Taispeánann duine acu do na \VUgras{póilíní} \textsuperscript{(64)} cá gcaithfear
\VUgras{íospartach} eile \textsuperscript{(65)} a iompar, fear dóiteáin cróga a thit
de dhréimire.

%%K t11-c2-03-1 p
3. --- Sonraíonn \VUgras{an t-oifigeach} \textsuperscript{(66)} an áit inar cheart an
\VUgras{caidéal gaile} \textsuperscript{(67)} atá ag teacht a shuíomh. Agus é ag
feitheamh leis an inneall cumhachtach sin, chuir sé \VUgras{caidéal láimhe}
\textsuperscript{(68)} i bhfearas, caidéal a scaoileann a \VUgras{phíobán} fada
\textsuperscript{(69)} tuilte uisce ar an tine. Oibríonn seisear fear dóiteáin é, fad
is a dhíríonn a gcomrádaí, a bhfuil an \VUgras{soc} \textsuperscript{(70)} ina lámh
aige, an \VUgras{scaird} \textsuperscript{(71)} ar lár an dóiteáin.

%%K t11-c2-04-1 p
4. --- Ar an taobh eile den amharclann cuireadh an \VUgras{dréimire} mór
\textsuperscript{(72)} in airde. Ligeann sé d'fhir an dóiteáin druidim leis na lasracha
a phléascann amach i measc guairneán \VUgras{deataigh} dhuibh \textsuperscript{(75)}.

%%K t11-tit-7 sub
\VUsaut{5.0mm}
\VUpk{10.2pt}{40}{Gaiscí Cróga.}
\VUsaut{3.4mm}

%%K t11-c2-05-1 p
5. --- Sa halla tosaigh a thosaigh \VUgras{an dóiteán} \textsuperscript{(74)}, halla
tosaigh na \VUgras{hamharclainne} \textsuperscript{(73)}, agus an \VUgras{ceoldráma} á
chleachtadh, ceoldráma a raibh a chéad \VUgras{léiriú} le bheith ann amárach. Ar an
dea-uair ní raibh sa halla ach beagán \VUgras{lucht féachana} \textsuperscript{(76)}.
Chuaigh na hainniseoirí ar foscadh ar an \VUgras{mbalcóin} \textsuperscript{(77)}, áit
a dtagann \VUgras{fir dhóiteáin} chróga \textsuperscript{(86)} i gcabhair orthu le
dréimire agus le téada.

%%K t11-c2-06-1 p
6. --- Faoin \VUgras{gcolúnáid} \textsuperscript{(78)}, áit a bhfógraíonn an
\VUgras{fógra} \textsuperscript{(79)} an léiriú, feictear \VUgras{ceoltóirí}
\textsuperscript{(81)} na \VUgras{ceolfhoirne} \textsuperscript{(80)} ag teitheadh,
agus a n-uirlisí á n-iompar acu: \VUgras{veidhlín} \textsuperscript{(81)},
\VUgras{fliúit} \textsuperscript{(82)}, \VUgras{dordveidhil} \textsuperscript{(83)},
\VUgras{trombón} \textsuperscript{(84)}, \VUgras{druma} \textsuperscript{(85)}, srl.
Déantar iarracht cuid den \VUgras{troscán} \textsuperscript{(87)} a shábháil.

%%K t11-c2-07-1 p
7. --- B'éigean do na \VUgras{haisteoirí} \textsuperscript{(89)} agus do na
\VUgras{ban-aisteoirí} \textsuperscript{(88)}, do na \VUgras{hamhránaithe} agus do na
\VUgras{mná amhránaíochta}, teitheadh agus a \VUgras{gculaith}
\textsuperscript{(90)} amharclainne orthu. Míníonn an teanór d'\VUgras{iriseoir}
\textsuperscript{(92)} conas a tharla an rud. Tháinig \VUgras{an tuairisceoir} ar
cosa in airde ón gcéad nóiméad i dteannta lucht údaráis: an \VUgras{préifeacht}
\textsuperscript{(91)}, an \VUgras{méara} \textsuperscript{(93)}, \VUgras{coimisinéir
na bpóilíní} \textsuperscript{(94)} agus \VUgras{oifigeach breithiúnach}
\textsuperscript{(95)}, a chuirfidh fiosrúchán ar bun chun na freagrachtaí a
shocrú.
